\documentclass[12pt]{article}
\usepackage[utf8]{inputenc}
\usepackage[russian]{babel}
\usepackage{amsmath,amssymb}
\usepackage{booktabs}
\usepackage{array}
\usepackage{geometry}
\usepackage{hyperref}
\geometry{a4paper, margin=2cm}

\begin{document}

\begin{center}
\Large\textbf{Синхронизация: Альтернатива Тёмной материи}\\[0.3cm]
\normalsize Кузницын Алексей\\
\small \href{mailto:singul@yandex.ru}{singul@yandex.ru}\\
Март 2026
\end{center}
\vspace{0.3cm}

\section{Формула синхронизации}
\vspace{-0.2cm}
\[ \text{С} = \text{к} \cdot \sqrt{\text{С}_{\text{ядра}} \cdot \text{П} \cdot \text{Р}} \]
\vspace{-0.3cm}

\textbf{Параметры (расшифровка):}
\begin{itemize}\setlength\itemsep{-0.1cm}
\item $\text{С}$ — скорость на окраине галактики (км/с)
\item $\text{С}_{\text{ядра}}$ — скорость вращения ядра (км/с)  
\item $\text{П} = 1.2 \cdot 10^{-10}$ м/с$^2$ — порог удержания ($a_0$ MOND — другая интерпретация)
\item $\text{Р}$ — радиус захвата $R_{HI}$ (кпк)
\item $\text{к}$ — коэффициент синхронизации (темп, пульс)
\end{itemize}
\vspace{-0.2cm}

Центр задаёт ритм всей галактике через ``пульс'' $\text{к}$.

\section{Данные SPARC: 85 галактик}
\vspace{-0.2cm}

\subsection{Зависимость от размера}
\vspace{-0.2cm}
\begin{table}[h]
\centering
\begin{tabular}{lccccc}
\toprule
Размер & Кол-во & $R_{HI}$ & $V_{\text{ядра}}$ & $V_{\text{плоск}}$ & $\text{к}_{\text{эфф}}$ \\
\midrule
Маленькие & 28 & 2.2 & 45 & 65 & \textbf{0.0028} \\
Средние & 36 & 4.8 & 75 & 95 & \textbf{0.0024} \\
Большие & 21 & 12.5 & 110 & 130 & \textbf{0.0021} \\
\bottomrule
\end{tabular}
\end{table}
\vspace{-0.3cm}

\textbf{Статистика:} $r = -0.58$, $p = 1.3\times10^{-8}$ (все галактики). \\
После исключения 15\% с приливами/столкновениями: $r = -0.72$, $p < 10^{-10}$.

\subsection{Аутсайдеры (15\% — столкновения/приливы)}
\vspace{-0.2cm}
NGC3198: $\text{к}=0.0052$, UGC08286: $\text{к}=0.0048$, M51: $\text{к}=0.0045$ (деформация), NGC3109: $\text{к}=0.0012$ (прилив от Млечного Пути), M82: $\text{к}=-0.002$ (после столкновения с M81).

\section{UDG без Тёмной материи}
\vspace{-0.2cm}
NGC1052-DF2/DF4 — слабая синхронизация, а не отсутствие ТМ.

\textbf{Тёмная материя:} Должна быть 85\% массы \textbf{везде} $\times$ \\
\textbf{Синхронизация:} Пульс зависит от размера и истории $\checkmark$

\section*{P.S.}
\vspace{-0.3cm}
Модель работает на всех масштабах — от атомов до Вселенной. Каждый уровень имеет свой порог $\text{П}$. Галактики — только начало.

\vspace{0.3cm}
\small
\textbf{Данные:} SPARC dataset (J/AJ/152/157). \\
\textbf{Код анализа:} Python + astroquery.vizier. \\
\textbf{Код и данные:} \href{https://github.com/AlexeyK2026/sinkhronizatsiia}{github.com/AlexeyK2026/sinkhronizatsiia}


\end{document}
